\chapter{Method and Progress}\label{ch:method}

This chapter describes what has been implemented and what has been completed by
the midterm stage. The main training pipeline now runs end-to-end, the selected
Vermont benchmark experiments have completed their planned training budgets, and
the evaluation utilities export unified test-month metrics, ranking tables, and
load-tracking figures. At this point, the project has moved beyond early
pipeline debugging: several baselines and communication variants can now be
compared within the same MAPPO framework using the completed Vermont result set.

\section{Problem Definition}

The CityLearn Annex 96 task is treated as a cooperative control problem with
$N=25$ building agents. At time step $t$, building $i$ observes its own local
state $o_i^t$ and chooses a continuous action $a_i^t$ for its controllable
devices. The actions of all buildings together affect battery states, indoor
conditions, and the total portfolio load. The goal is to make the total load
follow a daily reference profile $r_t$ without creating too many comfort
violations or other negative side effects.

The project objective is not only to get a high training reward. In practice,
the controller should reduce the testing-month tracking error, measured by NMBE
and CV-RMSE, while keeping comfort, ramping, peak demand, and fairness in a
reasonable range. Cost and carbon emissions are part of the Annex 96 secondary
metric set, but the CE1 schemas used here do not provide pricing or
carbon-intensity time series, so those two metrics are not interpreted in the
current experiments. This is why the method uses a centralized critic and
structured communication between actors.

\section{Environment and Baseline Pipeline}

The first completed component is the environment interface. CityLearn is wrapped
for several control and learning pipelines:

\begin{itemize}
  \item RBC notebook baseline for non-learning comparison;

For the current standard, grouped, and communication-aware MAPPO experiments,
the \texttt{mappo\_standard/env.py} wrapper returns local observations,
centralized observations, continuous action spaces, rewards, done masks, and
episode lengths in the tensor format expected by \texttt{on-policy-main}. The
training scripts then reuse \texttt{R\_MAPPO}, \texttt{R\_MAPPOPolicy}, and
\texttt{SharedReplayBuffer} from that implementation, instead of reimplementing
the PPO optimizer and buffer logic. This reduces implementation risk and keeps
the research focus on CityLearn adaptation, grouping, and communication.

\section{Grouped Actor Architecture}

The grouped actor design sits between two simple choices. If all buildings share
one actor, training is efficient, but the policy may ignore important
differences between buildings. If every building has its own actor, the policy
can be more flexible, but it needs more data and coordination becomes harder.
The grouped design shares an actor only within similar buildings.

Buildings are grouped with K-means using features related to flexibility, mainly
battery energy storage capacity and total HVAC capacity. After clustering,
buildings in group $G_k$ share the same actor network. All groups still share one
centralized critic. The critic sees the full portfolio state, while each actor
uses its local building observations and, when enabled, communication-enhanced
features.

This architecture has been implemented in the \texttt{mappo\_grouped} package.
The same grouping information is reused by communication-aware variants so that
communication effects can be compared against a consistent grouped baseline.

\section{MAPPO Learning Structure}

The learning structure follows centralized training and decentralized
execution. During training, the critic can see the full portfolio state and
learn whether the overall control decision is good. During execution, each
group actor only uses its own building observations, plus any communication
features, to choose continuous actions. After each episode, MAPPO uses the
collected experience to estimate which actions were better than expected, and
then updates the actors with PPO's conservative update rule.

The grouped replay buffer keeps the group-agent structure when samples are drawn
for PPO updates. This matters for communication: the communication layer needs
to know which agents were in the same group at the same time step. If the data is
flattened too early, that structure is lost. The buffer therefore keeps tensors
in a shape that supports both MAPPO updates and group communication.

\section{Baseline and Communication Variants}

The implemented methods are organized as a step-by-step experimental path. The
project first tests individual learning methods, then standard MAPPO, then
grouped MAPPO without communication, and finally several communication variants.
This order makes it easier to see what is added at each stage.

\subsection{Individual PPO and SAC Baselines}

The first learning baselines are independent PPO and independent SAC. In these
runs, each building learns without receiving information from other buildings.
This gives a simple reference point for what can be achieved without multi-agent
coordination.

\subsection{Standard MAPPO}

The next step is standard MAPPO. All 25 buildings reuse the same actor network
through parameter sharing, while the centralized critic sees the full portfolio
state during training. There is still no explicit communication module in this
baseline. It tests whether a basic MARL method already improves the task before
adding grouping or communication design.

\subsection{Grouped MAPPO Without Communication}

After standard MAPPO, buildings are divided into groups with K-means. Buildings
in the same group share one group actor, but the encoded feature is not changed
by any message passing:
\[
\tilde{h}_i = h_i .
\]
This tests whether grouping alone helps. It is also the direct no-communication
baseline for later communication modules.

In the communication variants below, the actor first encodes each building
observation into a hidden feature. The communication layer then modifies this
feature before the action head selects the final action:
\[
o_i \rightarrow \mathrm{encoder}_k \rightarrow h_i
\rightarrow \mathrm{communication} \rightarrow \tilde{h}_i
\rightarrow \mathrm{action\ head}_k \rightarrow a_i .
\]
Because the communication module is part of the actor network, it is trained
together with the actor.

\subsection{CommNet-Style Mean Communication}

For each group, the CommNet-style module lets each agent receive the average
hidden feature of the other agents:
\[
m_i = \frac{1}{|G_k|-1}\sum_{j \in G_k, j \ne i} h_j .
\]
This average message is passed through a small neural network and then added
back to the agent's own feature:
\[
\tilde{h}_i = h_i + f_\theta(m_i).
\]
This is the simplest communication baseline in the project.

\subsection{Weighted Same-Group and Other-Group Communication}

The project-specific communication module separates messages from similar
buildings and messages from the rest of the portfolio. For group $G_k$,
\[
m^{same}_k = \frac{1}{|G_k|}\sum_{j \in G_k} h_j,
\qquad
m^{other}_k = \frac{1}{N-|G_k|}\sum_{j \notin G_k} h_j .
\]
The combined message is
\[
m_k = \alpha m^{same}_k + \beta m^{other}_k,
\qquad \alpha > \beta \geq 0.
\]
The default implementation uses $\alpha=0.75$ and $\beta=0.25$. This means that
same-group information is weighted more strongly, while other-group information
is still kept as weaker global context.

\subsection{PowerNet-Style Structured Information Passing}

The PowerNet-style module tests a more structured way of sharing information. It
can use ring, chain, or full topologies to decide which other agents are included
in the message. For adjacency matrix $A$, the selected peer message is
\[
m_i = \frac{1}{\max(1,\sum_j A_{ij})}\sum_j A_{ij}g_\theta(h_j).
\]
The implementation then combines the agent's own feature with the peer message
using a learned residual update. The purpose is to avoid overwriting local
information while still allowing agents to coordinate.

\subsection{GAT-Style Graph Attention Communication}

The GAT-style module uses masked graph attention over a fixed similarity graph.
Inside each K-means group, buildings are fully connected, and a small number of
cross-group edges connect the closest groups. Each agent then learns attention
weights over its allowed neighbors instead of using a simple average. This makes
it possible to test whether graph-structured attention improves communication
selectivity in the CityLearn setting.

\subsection{TarMAC-Style Attention}

The TarMAC-style module uses attention. Each agent produces query, key, and
value features, and the attention weights are
\[
\alpha_{ij} =
\frac{\exp(q_i^\top k_j / \sqrt{d})}
{\sum_l \exp(q_i^\top k_l / \sqrt{d})},
\]
and the received context is
\[
c_i = \sum_j \alpha_{ij}v_j .
\]
The updated feature is computed from $h_i$ and $c_i$. Compared with the fixed
$\alpha$ and $\beta$ weights above, this module learns the communication weights
from the current features. It tests whether the agents can learn which other
agents are worth paying more attention to.

\subsection{DIAL-Style Differentiable Communication}

The DIAL-style module tests a different communication idea. Each agent first
turns its hidden feature into a message. During training, the message remains
continuous and differentiable, so gradients can pass through the communication
channel. During evaluation, the message can be discretized. In this project, the
DIAL module is adapted to MAPPO rather than implemented as the original
DQN-based DIAL algorithm. The purpose is to test whether trainable
communication signals help the grouped actors coordinate.

\section{Completed Implementation}

At the midterm stage, the following implementation work has been completed:

\begin{itemize}
  \item RBC notebook baseline for non-learning comparison;

The most important outcome so far is that the main methods can now be run and
compared through a common experimental interface. The completed Vermont
benchmark results show that communication changes the balance between tracking,
comfort, and secondary metrics. In the current selected benchmark set,
\texttt{mappo\_grouped\_comm\_v2\_vt\_500\_final} is the strongest overall
result, while TarMAC, weighted communication with
\mbox{$\alpha=0.90, \beta=0.10$}, and GAT also show competitive trade-offs.

\section{Current Limitations}

The current comparison is much more complete than the earlier debugging-stage runs,
but it is still not the last word on the full research question. The selected
Vermont experiments now have completed test-month exports, yet Texas RL
experiments and multi-seed validation are still future work. The CE1 schemas
used in the current pipeline also do not include pricing or grid-carbon time
series, so cost and carbon metrics cannot be interpreted without adding external
assumptions. Another limitation is that K-means currently uses static device
features; future work may replace or extend them with learned features from
building trajectories.
